\section{Ewolucja wag jako metoda trenowania agenta}
W rozdziale~\ref{chap:ea} opisano algorytmy ewolucyjne i ich warianty adaptacyjne. Kluczowym pytaniem jest jednak to, jak przekładają się one na konkretnego agenta będącego w stanie obsłużyć rozgrywkę w Hearthstone. Odpowiedź leży w doborze dwóch elementów: \textbf{reprezentacji agenta} oraz \textbf{schematu oceny}. Agent, którego optymalizacją zajmuje się niniejsza praca, jest opisywany przez wektor wag $\mathbf{w} \in \mathbb{R}^D$ przekazywany do ważonej funkcji oceny stanu planszy. W celu optymalizacji tych wag García-Sánchez i in. \cite{sanchez2020} stosują strategię ewolucyjną. Każdy kandydat jest oceniany przez serię symulowanych rozgrywek. Schemat tej procedury opisano w dalszej części pracy. Obejmuje on definicję cech wektora stanu, funkcję przystosowania oraz protokół koewolucji.

Celem rozdziału jest porównanie proponowanego podejścia z innymi metodami stosowanymi w literaturze. Omówiono podejście ewolucyjne będące bezpośrednią bazą eksperymentów, a następnie alternatywne podejścia zastosowane w ramach \textit{Hearthstone AI Competition} w edycjach 2018--2020: algorytm MCTS, algorytmy lookahead oraz inne wybrane rozwiązania. Zestawienie to pozwala wyraźnie odróżnić podejście \textit{offline} --- polegające na trenowaniu agenta przed meczem --- od podejść \textit{online}, które przeszukują drzewo gry w czasie podejmowania każdej decyzji. Pozwala to również uzasadnić wybór algorytmu ewolucyjnego jako metody badawczej.